\Subhead{1.4.3}{Усреднение в системах, содержащих быстрые и медленные движения}
\textbf{Принцип усреднения}: медленное движение определяется средним действием быстрых колебаний, но усреднять нужно с учётом вызываемого ими движения. Укороченные уравнения --- один пример этого принципа.

\paragraph{Медленное изменение параметров.}
Для $\ddot x+\omega^2(t)x=0$ при $\omega>0$, $|\dot\omega|\ll\omega^2$ подстановка почти гармонического движения и усреднение дают $\dot A/A=-\dot\omega/(2\omega)$. Поэтому
\begin{equation}
 x\simeq\frac C{\sqrt{\omega(t)}}\cos\left(\int^t\omega(t')\dd t'+\phi_0\right),
 \qquad J=E/\omega\simeq\mathrm{const}.
\end{equation}
Сохраняется действие, а энергия меняется из-за работы привода. В общем случае $J=(2\pi)^{-1}\oint p\dd q$. При сепаратрисе, $\omega\to0$ или резонансе это приближение нарушается.

\paragraph{Быстрая вибрация: маятник Капицы.}
Для $m\ddot x=-U'(x)+f(x)\cos\Omega t$ разделим $x=X+\xi$, где $X$ медленная, $\xi$ малая быстрая координата. Ведущий быстрый ответ $\xi=-f(X)\cos\Omega t/(m\Omega^2)$. Разложение $f(X+\xi)$ даёт средний член $\avg{f'\xi\cos\Omega t}=-f'f/(2m\Omega^2)$, откуда
\begin{equation}
 m\ddot X=-\frac{\dd U_{\rm eff}}{\dd X},\qquad
 U_{\rm eff}=U+\frac{f^2}{4m\Omega^2}.
\end{equation}
Для маятника с вертикальной вибрацией подвеса амплитуды $a$ и частоты $\Omega$, при угле $\vartheta$ от нижней вертикали,
\begin{equation}\label{eq:kapitza}
 U_{\rm eff}=-mgl\cos\vartheta+\frac{ma^2\Omega^2}{4}\sin^2\vartheta.
\end{equation}
Вторая производная у $\vartheta=\pi$ положительна при $a^2\Omega^2>2gl$: верхнее положение стабилизируется. Нужны $a/l\ll1$, $\Omega\gg\sqrt{g/l}$ и разделение времён. Нулевое среднее самой силы не означает нулевого среднего её действия.
